\documentclass[11pt,a4paper]{article}
\usepackage[utf8]{inputenc}
\usepackage[T1]{fontenc}
\usepackage[french]{babel}
\usepackage{lmodern}
\usepackage{amsmath,amssymb}
\usepackage[margin=2.5cm]{geometry}
\usepackage{microtype}
\usepackage{hyperref}
\usepackage{listings}
\lstset{basicstyle=\ttfamily\small, frame=single, breaklines=true,
  columns=fullflexible, keepspaces=true}
\title{Projet MLSF : résolution du jeu Meowdoku\\
  \large Question 3 : résolution et restitution d’une grille}
\author{}
\date{}

\begin{document}
\maketitle

\section{Du fichier texte à la solution}

La fonction \texttt{solve\_meowdoku} du fichier \texttt{meowdoku.py}
reçoit le chemin du fichier texte contenant la grille et celui du fichier
résultat. Elle réutilise les programmes des deux premières questions :
\begin{enumerate}
  \item \textbf{Lecture :} \texttt{read\_meowdoku} transforme le fichier
    en une matrice de régions. Le texte original est conservé pour être
    recopié dans le résultat.
  \item \textbf{Conversion et résolution :} \texttt{meowdoku\_to\_dimacs}
    écrit les contraintes dans un fichier DIMACS temporaire.
    \texttt{read\_dimacs} récupère les clauses, puis \texttt{dpll} cherche
    une valuation qui les satisfait. Si aucune solution n'existe,
    la fonction renvoie \texttt{None} sans écrire de résultat.
  \item \textbf{Décodage et vérification :} la valuation est complétée,
    puis les variables vraies sont converties en coordonnées de chats.
    Le programme vérifie les clauses SAT et les règles du jeu.
  \item \textbf{Restitution :} le fichier de sortie contient l'entrée
    originale, les coordonnées du chat de chaque région, puis une grille
    où \texttt{*} représente un chat et \texttt{.} une case vide.
    La fonction renvoie aussi le dictionnaire
    \texttt{\{région: (ligne, colonne)\}}.
\end{enumerate}

\section{Des variables vraies aux coordonnées des chats}

La partie centrale de cette question consiste à interpréter le modèle
renvoyé par le solveur. Chaque variable correspond à une case : si elle
vaut vrai, cette case contient un chat. Pour une grille de $m$ colonnes,
la numérotation définie à la question 1 est
\[
  v=(i-1)m+j,
\]
où $i$ et $j$ sont la ligne et la colonne, numérotées à partir de 1.
Pour retrouver ces coordonnées, on inverse cette relation :
\[
  i=\left\lfloor\frac{v-1}{m}\right\rfloor+1,
  \qquad j=((v-1)\bmod m)+1.
\]
Le code correspondant est :
\begin{lstlisting}[language=Python]
positions = [((v - 1) // m + 1, (v - 1) % m + 1)
             for v, value in model.items() if value]
\end{lstlisting}
\begin{itemize}
  \item \texttt{model.items()} parcourt les couples formés du numéro
    de variable \texttt{v} et de sa valeur booléenne \texttt{value}.
  \item \texttt{if value} conserve uniquement les variables vraies,
    donc les cases contenant un chat.
  \item \texttt{(v - 1) // m + 1} donne la ligne : la division entière
    compte les lignes complètes qui précèdent la case.
  \item \texttt{(v - 1) \% m + 1} donne la colonne : le reste de la
    division indique la position dans la ligne.
\end{itemize}
On soustrait 1 avant le calcul pour travailler à partir de zéro, puis on
ajoute 1 aux résultats pour obtenir des coordonnées commençant à 1.
Par exemple, avec $m=5$ colonnes et $v=8$ :
\[
  i=7\mathbin{//}5+1=2,\qquad j=7\bmod5+1=3.
\]
Si la variable 8 est vraie, un chat est donc placé en $(2,3)$.
La région de ce chat est ensuite obtenue par
\texttt{grid[i - 1][j - 1]}. Ces positions permettent d'écrire aussi
bien les coordonnées par région que les étoiles de la grille résolue.

\section{Utilisation}

Depuis le dossier du projet :
\begin{lstlisting}
python3 meowdoku.py grille.txt solution.txt --seed 0
\end{lstlisting}
L'option facultative \texttt{--seed} permet de reproduire les choix
aléatoires du solveur. Le programme écrit une solution vérifiée dans
\texttt{solution.txt} si la grille est satisfaisable.


\end{document}
